% Text C – António Guterres, Remarks to the UN Security Council on AI (18 July 2023), abridged
% Quelle: UN Press Release SG/SM/21880, https://press.un.org/en/2023/sgsm21880.doc.htm
% Übersetzen: xelatex text-c-guterres.tex (zweimal)
% ============================================================
% ab-vorlage.tex – gemeinsame Vorlage der Arbeitsblätter
% englisch/13/questions-on-the-text/arbeitsblaetter/
% Übersetzen mit:  xelatex <datei>.tex   (zweimal, wegen Seitenzahlen)
% Lösungsfassung:  Einzeiler-Datei <datei>-lsg.tex mit
%                  \def\mitloesung{}\input{<datei>}
% Schriften: Carlito (Fließtext), Poppins (Überschriften), TeX Gyre Pagella (Texte)
% ============================================================
\documentclass[10pt,a4paper]{article}
\usepackage[a4paper,top=12mm,bottom=13mm,left=14mm,right=14mm,headsep=4mm,headheight=12pt,footskip=7mm]{geometry}
\usepackage{fontspec}
\setmainfont{Carlito}
\newfontfamily\headfont{Poppins}[UprightFont=* Medium, BoldFont=* Medium]
\newfontfamily\headlight{Poppins Light}
\newfontfamily\lorafont{TeX Gyre Pagella}
\usepackage{xcolor}
\definecolor{prim}{HTML}{2563EB}
\definecolor{primdark}{HTML}{1E3A8A}
\definecolor{sec}{HTML}{7C3AED}
\definecolor{primtint}{HTML}{EFF6FF}
\definecolor{sectint}{HTML}{F5F3FF}
\definecolor{ink}{HTML}{1E293B}
\definecolor{muted}{HTML}{64748B}
\definecolor{rule}{HTML}{CBD5E1}
\definecolor{good}{HTML}{059669}
\definecolor{goodtint}{HTML}{ECFDF5}
\definecolor{warn}{HTML}{B45309}
\definecolor{warntint}{HTML}{FFFBEB}
\definecolor{bad}{HTML}{B91C1C}
\definecolor{badtint}{HTML}{FEF2F2}
\color{ink}
\usepackage[most]{tcolorbox}
\usepackage{tikz}
\usetikzlibrary{positioning,calc,arrows.meta,shapes.misc}
\usepackage{enumitem}
\setlist{nosep,leftmargin=1.1em,topsep=1pt}
\setlist[itemize]{label=\textcolor{prim}{\small\textbullet}}
\usepackage{tabularx,array,booktabs,colortbl}
\usepackage{multicol}
\usepackage{fancyhdr}
\usepackage{lastpage}
\usepackage{amssymb}
\usepackage{microtype}
\setlength{\parindent}{0pt}
\setlength{\parskip}{0pt}
\linespread{1.02}

% ---------- Kopf- und Fußzeile ----------
\newcommand{\wskopf}{English · Q13 · Science and Technology}
\pagestyle{fancy}
\fancyhf{}
\renewcommand{\headrulewidth}{0pt}
\fancyhead[L]{\footnotesize\headlight\textcolor{muted}{\wskopf}}
\fancyhead[R]{\footnotesize\headlight\textcolor{muted}{Name: \rule{38mm}{0.3pt}\quad Date: \rule{18mm}{0.3pt}}}
\fancyfoot[L]{\scriptsize\textcolor{muted}{edu-mrh.de/englisch/13/questions-on-the-text}}
\fancyfoot[R]{\scriptsize\textcolor{muted}{\thepage\,/\,\pageref*{LastPage}}}

% ---------- Titel ----------
% \wstitel{Kürzel}{Titel}{Untertitel}
\newcommand{\wstitel}[3]{%
  \begin{tcolorbox}[enhanced,colback=prim,colframe=prim,arc=2mm,boxrule=0pt,
     left=4mm,right=4mm,top=2.2mm,bottom=2.2mm,
     interior style={left color=primdark,right color=sec}]
    \color{white}{\headlight\footnotesize #1}\\[0.3mm]
    {\headfont\LARGE #2}\\[0.6mm]
    {\small #3}
  \end{tcolorbox}\vspace{1mm}}

% ---------- Boxen ----------
% \begin{wsbox}{Nummer}{Überschrift} ... \end{wsbox}
\newtcolorbox{wsbox}[3][]{enhanced,breakable,colback=white,colframe=rule,boxrule=0.5pt,arc=1.5mm,
  left=2.5mm,right=2.5mm,top=1.3mm,bottom=1.3mm,
  title={\headfont\normalsize\textcolor{white}{\colorbox{prim}{\,#2\,}}\;\textcolor{ink}{#3}},
  coltitle=ink,colbacktitle=white,attach title to upper={\par\vspace{0.6mm}},before skip=1.4mm,after skip=1.4mm,#1}
\newtcolorbox{tintbox}[1][]{enhanced,colback=primtint,colframe=primtint,boxrule=0pt,arc=1.2mm,
  left=2mm,right=2mm,top=1.2mm,bottom=1.2mm,#1}
\newtcolorbox{secbox}[1][]{enhanced,colback=sectint,colframe=sectint,boxrule=0pt,arc=1.2mm,
  left=2mm,right=2mm,top=1.2mm,bottom=1.2mm,#1}
\newtcolorbox{goodbox}[1][]{enhanced,colback=goodtint,colframe=goodtint,boxrule=0pt,arc=1.2mm,
  left=2mm,right=2mm,top=1.2mm,bottom=1.2mm,#1}
\newtcolorbox{warnbox}[1][]{enhanced,colback=warntint,colframe=warntint,boxrule=0pt,arc=1.2mm,
  left=2mm,right=2mm,top=1.2mm,bottom=1.2mm,#1}
\newtcolorbox{badbox}[1][]{enhanced,colback=badtint,colframe=badtint,boxrule=0pt,arc=1.2mm,
  left=2mm,right=2mm,top=1.2mm,bottom=1.2mm,#1}

\newcommand{\kl}[1]{{\headfont\small\textcolor{prim}{#1}}}      % kleine Zwischenüberschrift
\newcommand{\klsec}[1]{{\headfont\small\textcolor{sec}{#1}}}
\newcommand{\plus}{\textcolor{good}{\textbf{+}}\ }
\newcommand{\minus}{\textcolor{bad}{\textbf{–}}\ }
\newcommand{\pfeil}{\textcolor{prim}{$\rightarrow$}\ }
\newcommand{\linie}[1][\linewidth]{\rule{#1}{0.3pt}}
% ---------- Lösungsschalter ----------
% \lsg[Breite]{Text}        Lücke mit Unterstrich; Lösung erscheint nur mit \mitloesung
% \lsglinien{n}{Text}       n Schreiblinien; Lösung wird auf die Linien geschrieben
% Beide reservieren in beiden Fassungen exakt denselben Platz.
\newif\ifloesung
\ifdefined\mitloesung \loesungtrue \else \loesungfalse \fi
\newcommand{\lsgstil}{\color{sec}\itshape}
\newcommand{\lsg}[2][30mm]{%
  \makebox[#1][l]{\rlap{\raisebox{-0.8mm}{\textcolor{rule}{\rule{#1}{0.4pt}}}}%
  \ifloesung{\lsgstil\small #2}\fi}}
\newlength{\lsgzeile}\setlength{\lsgzeile}{7mm}
\newcommand{\lsglinien}[2]{%
  \par\noindent
  \begin{tikzpicture}[x=1mm,y=1mm]
    \pgfmathsetmacro{\hoehe}{#1*7}
    \path[use as bounding box] (0,0.5) rectangle (\linewidth,-\hoehe);
    \foreach \i in {1,...,#1} {\draw[rule,line width=0.4pt] (0,-\i*7) -- (\linewidth,-\i*7);}
    \ifloesung
      \node[anchor=north west,inner sep=0pt,text width=\linewidth,align=left,
            font=\lsgstil\small] at (0,-3.3) {\setlength{\baselineskip}{\lsgzeile}#2\par};
    \fi
  \end{tikzpicture}\par}
\newcommand{\schreiblinien}[1]{\lsglinien{#1}{}}

\geometry{left=22mm,right=58mm,marginparwidth=40mm,marginparsep=6mm}
\usepackage[modulo]{lineno}
\renewcommand{\linenumberfont}{\normalfont\scriptsize\color{muted}}
\setlength{\linenumbersep}{5mm}
\newcommand{\glosse}[1]{\marginpar{\raggedright\footnotesize\color{muted}\setlength{\baselineskip}{9pt}#1}}
\fancyhfoffset[R]{44mm}
\begin{document}
\begin{tcolorbox}[enhanced,colback=primtint,colframe=primtint,arc=1.5mm,left=3mm,right=3mm,top=2mm,bottom=2mm,
   width=\linewidth+46mm]
{\headlight\footnotesize\textcolor{prim}{Text C · Speech · Choice text for “Questions on the text”}}\\[0.6mm]
{\headfont\large Secretary-General's remarks to the Security Council on Artificial Intelligence}\\[1mm]
{\small On 18 July 2023, the UN Security Council – the body responsible for international peace and security – held its first-ever debate on artificial intelligence. UN Secretary-General António Guterres opened the meeting with this speech.}
\end{tcolorbox}
\vspace{2mm}
{\lorafont\fontsize{10.4}{13.9}\selectfont
\setlength{\parindent}{1.2em}
\linenumbers
I thank the United Kingdom for convening\glosse{\textbf{to convene} to call a formal meeting} the first debate on artificial intelligence (AI) ever held in this Council.\par
I have been following the development of AI for some time. Indeed, I told the General Assembly six years ago that AI would have a dramatic impact on sustainable development, the world of work and the social fabric\glosse{\textbf{social fabric} the structure of relationships that holds a society together}. But, like everyone here, I have been shocked and impressed by the newest form of AI, generative AI, which is a radical advance in its capabilities.\par
The speed and reach of this new technology in all its forms are utterly unprecedented\glosse{\textbf{unprecedented} never having happened before}. It has been compared to the introduction of the printing press. But, while it took more than 50 years for printed books to become widely available across Europe, ChatGPT reached 100 million users in just two months.\par
The finance industry estimates AI could contribute between \$10 and \$15 trillion to the global economy by 2030. Almost every Government, large company and organization in the world is working on an AI strategy.\par
But, even its own designers have no idea where their stunning technological breakthrough may lead. It is clear that AI will have an impact on every area of our lives, including the three pillars of the United Nations. It has the potential to turbocharge\glosse{\textbf{to turbocharge} to make much faster or stronger} global development, from monitoring the climate crisis to breakthroughs in medical research. It offers new potential to realize human rights, particularly to health and education. But, the High Commissioner for Human Rights has expressed alarm over evidence that AI can amplify bias\glosse{\textbf{to amplify} to make stronger}, reinforce discrimination and enable new levels of authoritarian surveillance.\par
Today’s debate is an opportunity to consider the impact of artificial intelligence on peace and security — where it is already raising political, legal, ethical and humanitarian concerns. I urge the Council to approach this technology with a sense of urgency, a global lens, and a learner’s mindset\glosse{\textbf{mindset} attitude, way of thinking}. Because what we have seen is just the beginning. Never again will technological innovation move as slow as it is moving today.\par
AI is being put to work in connection with peace and security, including by the United Nations. It is increasingly being used to identify patterns of violence, monitor ceasefires\glosse{\textbf{ceasefire} agreement to stop fighting} and more, helping to strengthen our peacekeeping, mediation and humanitarian efforts.\par
But, AI tools can also be used by those with malicious intent\glosse{\textbf{malicious} meaning to do harm}. AI models can help people to harm themselves and each other, at massive scale. Let’s be clear: the malicious use of AI systems for terrorist, criminal or State purposes could cause horrific levels of death and destruction, widespread trauma and deep psychological damage on an unimaginable scale.\par
AI-enabled cyberattacks are already targeting critical infrastructure\glosse{\textbf{critical infrastructure} vital systems, e.g. power, water, hospitals} and our own peacekeeping and humanitarian operations, causing great human suffering. The technical and financial barriers to access are low, including for criminals and terrorists. Both military and non-military applications of AI could have very serious consequences for global peace and security.\par
The advent\glosse{\textbf{advent} arrival, beginning} of generative AI could be a defining moment for disinformation and hate speech — undermining truth, facts and safety, adding a new dimension to the manipulation of human behaviour and contributing to polarization and instability on a vast scale.\par
Deepfakes are just one new AI-enabled tool that, if unchecked, could have serious implications for peace and stability. And the unforeseen consequences of some AI-enabled systems could create security risks by accident. Look no further than social media. Tools and platforms that were designed to enhance human connection are now used to undermine elections, spread conspiracy theories and incite hatred\glosse{\textbf{to incite} to encourage violence or hate} and violence.\par
Malfunctioning AI systems are another huge area of concern. And the interaction between AI and nuclear weapons, biotechnology, neurotechnology and robotics is deeply alarming. Generative AI has enormous potential for good and evil at scale. Its creators themselves have warned that much bigger, potentially catastrophic and existential risks lie ahead. Without action to address these risks, we are derelict\glosse{\textbf{to be derelict in sth} to fail in one’s duty} in our responsibilities to present and future generations.\par
The international community has a long history of responding to new technologies with the potential to disrupt our societies and economies. We have come together at the United Nations to set new international rules, sign new treaties and establish new global agencies. While many countries have called for different measures and initiatives around the governance of AI, this requires a universal approach.\par
And questions of governance will be complex for several reasons. First, powerful AI models are already widely available to the general public. Second, unlike nuclear material and chemical and biological agents, AI tools can be moved around the world leaving very little trace. And third, the private sector’s leading role in AI has few parallels in other strategic technologies.\par
[…] The best approach would address existing challenges while also creating the capacity to monitor and respond to future risks. It should be flexible and adaptable, and consider technical, social and legal questions. It should integrate the private sector, civil society, independent scientists and all those driving AI innovation.\par
The need for global standards and approaches makes the United Nations the ideal place for this to happen. The Charter of the United Nations’ emphasis on protecting succeeding generations gives us a clear mandate to bring all stakeholders together around the collective mitigation\glosse{\textbf{mitigation} making sth less harmful} of long-term global risks. AI poses just such a risk.\par
I therefore welcome calls from some Member States for the creation of a new United Nations entity\glosse{\textbf{entity} organisation, body} to support collective efforts to govern this extraordinary technology, inspired by such models as the International Atomic Energy Agency (IAEA), the International Civil Aviation Organization (ICAO) or the Intergovernmental Panel on Climate Change.\par
[…] Let’s be honest: There is a huge skills gap around AI in Governments and other administrative and security structures that must be addressed at the national and global levels. A new United Nations entity would gather expertise and put it at the disposal of the international community. And it could support collaboration on the research and development of AI tools to accelerate sustainable development.\par
[…] I urge this Council to exercise leadership on artificial intelligence and show the way towards common measures for the transparency, accountability, and oversight of AI systems. We must work together for AI that bridges social, digital and economic divides, not one that pushes us further apart.\par
I urge you to join forces and build trust for peace and security. We need a race to develop AI for good: to develop AI that is reliable and safe and that can end poverty, banish hunger\glosse{\textbf{to banish} to get rid of}, cure cancer and supercharge climate action [and] an AI that propels us towards the Sustainable Development Goals\glosse{\textbf{Sustainable Development Goals} 17 UN goals for 2030, e.g.\ no poverty}. That is the race we need, and that is a race that is possible and achievable. Thank you.\par
\nolinenumbers}
\vspace{1.5mm}
\hfill{\small António Guterres, remarks to the UN Security Council, New York, 18 July 2023\\
\hspace*{\fill}Source: UN press release SG/SM/21880 (abridged)}
\end{document}
